\scanpage{92}
drais l'interpréter comme l'extension \uline{universelle} (de Serre) de $B$ par un fibré vectoriel.

Notons d'abord que le schéma $B^{*}$ des translations de $A$ opère \unsure{triv\textsuperscript{t}|trivialt|trivialement} sur l'extension $\deleted{\unsure{M_{A/S}|\underline E_{A/S}}}\inserted{\mathcal M_{A/S}^{0}\deleted{\unsure{=|\underline E}}}$ (car il\deleted{s} opère \unsure{triv\textsuperscript{t}|trivialement} sur $\underline{\mathrm{Pic}}^{0}_{A/S}$ et sur $\underline{\omega}_{A/S}$\deleted{)}, et $\unsure{\forall|\mathrm{tt}}$ hom de $B$ dans $\underline{\omega}_{X/S}$ est trivial). Donc par théorie de descente on voit que $\mathcal M^{0}_{A/S}$ est can isom à $\mathcal M^{0}_{B^{*}/S}$, donc c'est un invariant qui ne dépend que du \deleted{\unsure{le}} schéma abélien $B$ sur $S$.

Pour calculer la classe de l'extension obtenue de $B$ par $\underline{\omega}_{A/S}$ comme élément de $H^{0}(S,\underbrace{R^{1}f_{B*}^{\deleted{\unsure{1}}}(\underline{\omega}_{A/S}\otimes_{\mathcal O_S}\mathcal O_B)}_{\unsure{(R^{1}f_{B*}\mathcal O_{B})|R^{1}f_{B*}(\mathcal O_{B})|R^{1}f_{B*}(\underline{\mathcal O}_{B})}^{\vee}\otimes R^{1}f_{B*}(\mathcal O_{B})})$, on regarde \inserted{par la recette générale} le faisceau inversible \unsure{canonique|unique|de Poincaré} sur $P=A\times_{S}B$, et on écrit
\diagram{\begin{tikzcd}[ampersand replacement=\&,column sep=small,row sep=small]
A \& \& B \& A\arrow[d,equal] \&[-0.8em] \times_{S}B\\
\& S\arrow[ul,"f_A"']\arrow[ur,"f_B",dash] \& \& B^{*} \& {}
\end{tikzcd}}
l'obstruction à mettre dessus une connexion relativement à $B$, i.e. la classe de cohomologie qui lui correspond dans $H^{1}(P,\underline{\Omega}^{1}_{P/B})$ [puis on localise sur $S$].

Calculons donc la classe de \unsure{cohomologie} de $\underline L$ dans\deleted{\unsure{à}} $H^{1}(A\times_{S}B,\underline{\Omega}^{1}_{A\times_{S}B/S})$\deleted{$\unsure{\in H|\text{de }L}$}, localisé en $S$.

% archive sheet number: 91

